%
% Only the most commonly-used macros. Needed by everything else.
%
\ifx\havemjocommon\undefined
\def\havemjocommon{1}

\ifx\mathbb\undefined
  \usepackage{amsfonts}
\fi

\ifx\restriction\undefined
  \usepackage{amssymb}
\fi

% Place the argument in matching left/right parentheses.
\newcommand*{\of}[1]{ \left({#1}\right) }

% Group terms using parentheses.
\newcommand*{\qty}[1]{ \left({#1}\right) }

% Group terms using square brackets.
\newcommand*{\sqty}[1]{ \left[{#1}\right] }

% A pair of things.
\newcommand*{\pair}[2]{ \left({#1},{#2}\right) }

% A triple of things.
\newcommand*{\triple}[3]{ \left({#1},{#2},{#3}\right) }

% A four-tuple of things.
\newcommand*{\quadruple}[4]{ \left({#1},{#2},{#3},{#4}\right) }

% A five-tuple of things.
\newcommand*{\quintuple}[5]{ \left({#1},{#2},{#3},{#4},{#5}\right) }

% A six-tuple of things.
\newcommand*{\sextuple}[6]{ \left({#1},{#2},{#3},{#4},{#5},{#6}\right) }

% A seven-tuple of things.
\newcommand*{\septuple}[7]{ \left({#1},{#2},{#3},{#4},{#5},{#6},{#7}\right) }

% A free-form tuple of things. Useful for when the exact number is not
% known, such as when \ldots will be stuck in the middle of the list,
% and when you don't want to think in column-vector terms, e.g. with
% elements of an abstract Cartesian product space.
\newcommand*{\tuple}[1]{ \left({#1}\right) }

% The "least common multiple of" function. Takes a nonempty set of
% things that can be multiplied and ordered as its argument. Name
% chosen for synergy with \gcd, which *does* exist already.
\newcommand*{\lcm}[1]{ \operatorname{lcm}\of{{#1}} }
\ifdefined\newglossaryentry
  \newglossaryentry{lcm}{
    name={\ensuremath{\lcm{X}}},
    description={the least common multiple of the elements of $X$},
    sort=l
  }
\fi

% The factorial operator.
\newcommand*{\factorial}[1]{ {#1}! }

% Restrict the first argument (a function) to the second argument (a
% subset of that functions domain). Abused for polynomials to specify
% an associated function with a particular domain (also its codomain,
% in the case of univariate polynomials).
\newcommand*{\restrict}[2]{{#1}{\restriction}_{#2}}
\ifdefined\newglossaryentry
  \newglossaryentry{restriction}{
    name={\ensuremath{\restrict{f}{X}}},
    description={the restriction of $f$ to $X$},
    sort=r
  }
\fi

%
% Product spaces
%
% All of the product spaces (for example, R^n) that follow default to
% an exponent of ``n'', but that exponent can be changed by providing
% it as an optional argument. If the exponent given is ``1'', then it
% will be omitted entirely.
%

% The natural n-space, N x N x N x ... x N.
\newcommand*{\Nn}[1][n]{
  \mathbb{N}\if\detokenize{#1}\detokenize{1}{}\else^{#1}\fi
}

\ifdefined\newglossaryentry
  \newglossaryentry{N}{
    name={\ensuremath{\Nn[1]}},
    description={the set of natural numbers},
    sort=N
  }
\fi

% The integral n-space, Z x Z x Z x ... x Z.
\newcommand*{\Zn}[1][n]{
  \mathbb{Z}\if\detokenize{#1}\detokenize{1}{}\else^{#1}\fi
}

\ifdefined\newglossaryentry
  \newglossaryentry{Z}{
    name={\ensuremath{\Zn[1]}},
    description={the ring of integers},
    sort=Z
  }
\fi

% The rational n-space, Q x Q x Q x ... x Q.
\newcommand*{\Qn}[1][n]{
  \mathbb{Q}\if\detokenize{#1}\detokenize{1}{}\else^{#1}\fi
}

\ifdefined\newglossaryentry
  \newglossaryentry{Q}{
    name={\ensuremath{\Qn[1]}},
    description={the field of rational numbers},
    sort=Q
  }
\fi

% The real n-space, R x R x R x ... x R.
\newcommand*{\Rn}[1][n]{
  \mathbb{R}\if\detokenize{#1}\detokenize{1}{}\else^{#1}\fi
}

\ifdefined\newglossaryentry
  \newglossaryentry{R}{
    name={\ensuremath{\Rn[1]}},
    description={the field of real numbers},
    sort=R
  }
\fi


% The complex n-space, C x C x C x ... x C.
\newcommand*{\Cn}[1][n]{
  \mathbb{C}\if\detokenize{#1}\detokenize{1}{}\else^{#1}\fi
}

\ifdefined\newglossaryentry
  \newglossaryentry{C}{
    name={\ensuremath{\Cn[1]}},
    description={the field of complex numbers},
    sort=C
  }
\fi

% The n-dimensional product space of a generic field F.
\newcommand*{\Fn}[1][n]{
  \mathbb{F}\if\detokenize{#1}\detokenize{1}{}\else^{#1}\fi
}

\ifdefined\newglossaryentry
  \newglossaryentry{F}{
    name={\ensuremath{\Fn[1]}},
    description={a generic field},
    sort=F
  }
\fi


% An indexed arbitrary binary operation such as the union or
% intersection of an infinite number of sets. The first argument is
% the operator symbol to use, such as \cup for a union. The second
% argument is the lower index, for example k=1. The third argument is
% the upper index, such as \infty. Finally the fourth argument should
% contain the things (e.g. indexed sets) to be operated on.
\newcommand*{\binopmany}[4]{
  \mathchoice{ \underset{#2}{\overset{#3}{#1}}{#4} }
             { {#1}_{#2}^{#3}{#4} }
             { {#1}_{#2}^{#3}{#4} }
             { {#1}_{#2}^{#3}{#4} }
}


% The four standard (UNLESS YOU'RE FRENCH) types of intervals along
% the real line.
\newcommand*{\intervaloo}[2]{ \left({#1},{#2}\right) } % open-open
\newcommand*{\intervaloc}[2]{ \left({#1},{#2}\right] } % open-closed
\newcommand*{\intervalco}[2]{ \left[{#1},{#2}\right) } % closed-open
\newcommand*{\intervalcc}[2]{ \left[{#1},{#2}\right] } % closed-closed


\fi
